% ============================================================================
% Section 1 - Proposed Approach
% Packages used: amsmath, booktabs, tabularx, url (for \path), natbib (for \citet)
% ============================================================================
\section{Proposed Approach}
\label{sec:approach}

\subsection{Task input and output}
\label{sec:approach-io}

For each target word $w$ the input is the set of its usages
$U_w=\{u_1,\dots,u_n\}$. A usage $u_i=(s_i,a_i,b_i,\pi_i)$ is a sentence $s_i$,
the character span $[a_i,b_i)$ of the target in it, and its period
$\pi_i\in\{\pi_1,\dots,\pi_{10}\}$. The ten periods are five four-year windows
between 1880 and 1915 and five two-year windows between 2006 and 2023. For
Subtask~2 the input also includes one sense definition $\delta_w$. The system
produces three outputs, matching what the task scores:
\begin{itemize}
  \item \textbf{Subtask 1a:} a sense label $c_i$ for every usage. The labels are
        induced cluster identifiers, not dictionary senses.
  \item \textbf{Subtask 1b:} a change score for every pair of consecutive periods.
        The official scorer computes the Jensen--Shannon divergence (JSD)
        between the sense distributions of two periods directly from the 1a
        labels, so the 1b prediction is fully determined by the clustering. We
        also write a clustering-free graded score, the average pairwise distance
        (APD), as a local diagnostic.
  \item \textbf{Subtask 2:} a label $\hat y_i\in\{0,1\}$ for every usage that says
        whether $u_i$ expresses the sense $\delta_w$.
\end{itemize}

\subsection{Pipeline}
\label{sec:approach-pipeline}

The system follows the four-step sense-based pipeline described by
\citet{periti2024towards}: extract usages, embed them, cluster them, and compare
periods. What sets it apart from the pretrained baseline is a
\emph{fine-tuned usage encoder} in place of a frozen one
(Section~\ref{sec:architecture}, Figure~\ref{fig:architecture}).

\paragraph{1. Preprocessing.} Every system reads the data through the shared
loader. It repairs five misaligned target offsets, maps the five pre-1906
Swedish spelling patterns (\textit{hv-}, \textit{-fv-}, \textit{-dt},
\textit{af}, \textit{qv-}) to their modern forms while keeping the target offsets
valid, and keeps the original case. The target is then wrapped in two special
tokens, \texttt{<t>}~target~\texttt{</t>}, so that the encoder knows which word
to represent \citep{cassotti2023xllexeme}.

\paragraph{2. Embedding extraction.} An XLM-R base encoder
\citep{conneau2020xlmr} encodes each usage independently, as a bi-encoder
\citep{reimers2019sbert}. The usage vector $\mathbf{v}_i\in\mathbb{R}^{768}$ is
the mean of the last-layer hidden states of the sub-words that overlap the
target span. The definition $\delta_w$ is encoded the same way as a
pseudo-usage ``$w$: $\delta_w$'', with $w$ as the target.

\paragraph{3. Sense clustering.} The usages of a word from \emph{all} periods are
clustered together with average-linkage agglomerative clustering on cosine
distance. Clustering every period separately would give cluster identifiers that
cannot be compared across periods, the alignment problem described by
\citet{periti2024towards}. The dendrogram is cut at distance $t$. Clusters with
fewer than $m$ usages are dissolved, and each of their usages joins the large
cluster with the most similar centroid. A word can therefore end up with a single
sense. The cut $t$ and the minimum size $m$ are chosen per fold on the validation
word (Section~\ref{sec:approach-selection}).

\paragraph{4. Sense distribution per period.} For period $\pi$, the sense
distribution $P_\pi$ is the share of the period's usages in each cluster.

\paragraph{5. Change scoring.} For consecutive periods $\pi_k,\pi_{k+1}$:
\begin{align}
  \mathrm{JSD}(P_{\pi_k},P_{\pi_{k+1}}) &= \tfrac12 D_{\mathrm{KL}}(P_{\pi_k}\,\|\,M)
     + \tfrac12 D_{\mathrm{KL}}(P_{\pi_{k+1}}\,\|\,M),\quad M=\tfrac12(P_{\pi_k}+P_{\pi_{k+1}}),
     \label{eq:jsd}\\
  \mathrm{APD}(\pi_k,\pi_{k+1}) &= \frac{1}{|U_{\pi_k}|\,|U_{\pi_{k+1}}|}
     \sum_{u\in U_{\pi_k}}\sum_{u'\in U_{\pi_{k+1}}}\bigl(1-\cos(\mathbf v_u,\mathbf v_{u'})\bigr).
     \label{eq:apd}
\end{align}
JSD (base~2) is the quantity the official 1b scorer ranks
\citep{lin1991divergence,giulianelli2020analysing}. APD needs no clustering
\citep{kutuzov2020uio}. Comparing the two shows whether a 1b error comes from the
representation or from the clustering step.

\paragraph{6. Subtask 2.} A usage is labelled positive if
$\cos(\mathbf v_i,\mathbf v_{\delta_w})\ge\tau$. The threshold $\tau$ is a single
absolute value tuned on the nine non-test words of the fold to maximise the
official macro F1. It is not a within-word rank, so the share of positive
predictions can differ between words.

\paragraph{7. Prediction.} The labels are written in the official file format and
scored with the shared evaluation script.

\subsection{Training objective}
\label{sec:approach-training}

The encoder is trained as a bi-encoder on usage pairs of the same word, following
XL-LEXEME \citep{cassotti2023xllexeme}, which learns from Word-in-Context
pairs \citep{pilehvar2019wic}. For a pair $(u,u')$ with gold sense sets $S_u$ and
$S_{u'}$ we use a \emph{soft} pair label and the contrastive loss of
\citet{hadsell2006contrastive} on cosine distance:
\begin{equation}
  y=\frac{|S_u\cap S_{u'}|}{|S_u\cup S_{u'}|},\qquad
  d=1-\cos(\mathbf v_u,\mathbf v_{u'}),\qquad
  \mathcal L = y\,d^{2} + (1-y)\,\max(0,\,m_c-d)^{2},
  \label{eq:loss}
\end{equation}
with margin $m_c=0.5$. Usages that share a sense are pulled together, and usages
with no shared sense are pushed at least $m_c$ apart. Of the development usages,
13.2\% carry two or three senses. For these, $y$ lies between 0 and 1, so such a
pair is neither forced together nor apart. With hard labels ($y\in\{0,1\}$),
Eq.~\eqref{eq:loss} is the loss used by XL-LEXEME.

Training has two stages.
\begin{itemize}
  \item \textbf{Stage A (synthetic English).} The organisers' optional synthetic
        English set \citep{tahmasebi2026semeval} contains 672{,}420 generated
        sentences with a sense identifier for every sentence. We drop the 5.3\%
        of lines whose marked span is not the target word (e.g.\ \textit{asked}
        marked for \textit{question}), keep words with at least two senses, and
        keep at most 20 sentences per sense. This leaves 2{,}193 training words
        (172{,}132 sentences), plus 68 held-out words that are only used to
        monitor this stage. We train on 50{,}000 sampled pairs for one epoch.
        Stage A uses no Swedish data, so it is run once per seed and shared by
        all folds.
  \item \textbf{Stage B (Swedish, per fold).} Starting from the Stage A weights,
        the encoder is fine-tuned on pairs from the fold's eight training words
        only, with 4{,}000 freshly sampled pairs per epoch and at most four
        epochs. After every epoch, and once before the first, we compute the AUC
        with which cosine similarity separates same-sense from different-sense
        pairs of the \emph{validation} word. Training stops after two epochs
        without improvement, and the best weights are kept. If no epoch improves
        on the Stage A weights, those weights are kept.
\end{itemize}
Pairs are drawn within one word: the word is chosen uniformly, then an anchor
usage, then a partner, with half of the pairs positive. Two usages with the same
sentence are never paired, and sentences whose duplicate copies received
conflicting gold labels are excluded. We use AdamW
\citep{loshchilov2019adamw}, batches of 16 pairs, linear warm-up over 10\% of the
steps, gradient clipping at 1.0 and mixed precision. The 192M-parameter word-piece
embedding matrix stays frozen, and the transformer layers (about 86M parameters)
are updated. Freezing the embeddings keeps the representation of the vocabulary
of the other five test languages unchanged and reduces memory use.

\paragraph{Era-aware sampling.} To stop the encoder from using era as a cue for
``same sense'' (EDA F4, F8), a positive partner is drawn from the other era with
probability 0.7 when one exists, and a negative partner from the same era with
probability 0.7. The eras are before and after 1950 for Swedish and before and
after 1900 for the synthetic set.

\subsection{Choices made without the test word}
\label{sec:approach-selection}

Each fold has one test word, one validation word and eight training words
(Section~7). The cut $t$ is the distance at which the pairwise decision
``same sense iff $d\le t$'' has the highest balanced accuracy on the validation
word's usage pairs. It is taken from the 2nd to the 98th percentile of the
observed distances. The minimum cluster size $m\in\{5,10,15,20,30\}$ is then
chosen by the validation word's BCubed F1, with ties going to the larger $m$.
The Subtask~2 threshold $\tau$ is tuned on the nine non-test words. The test word
is used only once, for scoring. Three alternative selection rules are
implemented for the ablations: the baseline procedure on the nine non-test words,
selection by validation BCubed F1, and a fixed cut $t=m_c/2$.

\subsection{Handling the challenges of the data}
\label{sec:approach-challenges}

\paragraph{Class and sense imbalance.} Most words have one dominant sense: one
cluster per word already scores 0.756 BCubed F1 (EDA F6). The minimum cluster
size lets the method return a single sense when no well-separated second sense
exists, and the cut is chosen per word rather than fixed. For Subtask~2 the
target-sense share ranges from 5.7\% to 94.6\% across words (EDA F7). The absolute
threshold $\tau$ lets the predicted positive rate follow the word. It is tuned on
the official macro F1 over both classes, which no longer rewards predicting
everything positive (that now scores 0.236). During training, words are sampled
uniformly and pairs are balanced 50/50, so large words and dominant senses do
not dominate the loss.

\paragraph{Old spellings and noisy historical text.} Pre-1906 spellings are
normalised before encoding (EDA F4). Sub-word tokenisation keeps OCR-damaged
words usable, since an unseen form is split into known pieces instead of becoming
an unknown token. Era-aware sampling removes era as a shortcut during training.

\paragraph{Long contexts.} We encode the whole sentence, because the encoder is
trained on full sentences, and truncate at 256 sub-words. No development sentence
is truncated: the longest has 571 characters. Texts longer than 1{,}000
characters are cropped around the target before tokenisation, so the target can
never be cut off. A fixed $\pm k$-word window is available in the configuration
(\texttt{context\_window}) for the ablation of EDA F3, where historical sentences
are 39\% longer than modern ones.

\paragraph{Different numbers of usages per period.} The organisers capped the
development data at 30 usages per word and period, and some periods have fewer
than 10 (EDA F1). Clustering all periods together uses every usage to estimate
the senses, so a small period does not get its own unstable clustering. JSD
compares relative frequencies, so period sizes do not enter the score directly.
The official scorer ignores periods with fewer than 10 usages. A training-only
cap (\texttt{max\_usages\_per\_period}) exists but is not needed at 30. It is never
applied to the prediction set, because the scorer requires a label for every
usage.

\subsection{Preprocessing decisions}
\label{sec:approach-preprocessing}

Table~\ref{tab:preprocessing} lists the preprocessing settings. All of them are
in the shared configuration file (\path{configs/default.yaml}) and are the same
for every system.

\begin{table}[ht]
\centering
\small
\begin{tabularx}{\linewidth}{@{}lXl@{}}
\toprule
\textbf{Decision} & \textbf{Setting and reason} & \textbf{Config key}\\
\midrule
Lowercasing & No. Assignment~1 lower-cased together with spelling normalisation and could not separate the two effects; case also marks proper nouns (\textit{Fröken Johanna}). & \texttt{lowercase}\\
Spelling normalisation & Yes, five pre-1906 patterns (EDA F4: archaic spellings occur in 71.5\% of historical usages and 0.8\% of modern ones). & \texttt{spelling\_normalisation}\\
Context window & Whole sentence, at most 256 sub-words; target never truncated. & \texttt{context\_window}\\
Max.\ usages per period & None for prediction (every usage must be labelled); training cap available but not needed (data capped at 30). & \texttt{max\_usages\_per\_period}\\
Target marking & \texttt{<t>} \dots\ \texttt{</t>} around the repaired target span. & \texttt{target\_markers}\\
Synthetic data & At most 20 sentences per sense; lines whose span is not the target are removed. & \texttt{synthetic.max\_per\_sense}\\
\bottomrule
\end{tabularx}
\caption{Preprocessing decisions of the proposed system.}
\label{tab:preprocessing}
\end{table}

\subsection{Expected limitations}
\label{sec:approach-limitations}

\begin{itemize}
  \item \textbf{Rare senses are merged.} A sense with fewer than $m$ usages
        (at least 5) cannot form its own cluster. \textit{suga} has 12 gold senses
        in 129 usages, so most of them will be merged.
  \item \textbf{Modern training text.} XLM-R was pretrained on modern web text,
        and the synthetic set is modern English. Historical Swedish, including
        OCR errors that spelling normalisation does not fix, is far from both.
  \item \textbf{Few labelled words.} Stage B sees only eight Swedish words per
        fold, and early stopping and clustering choices rest on a single
        validation word. These choices are therefore noisy, and results can vary
        between seeds.
  \item \textbf{One sense per usage.} The clustering assigns one cluster per usage,
        although 13.2\% of usages have several gold senses.
  \item \textbf{Subtask~2 is only indirectly trained.} The definition vector comes
        from the same encoder, but no training pair contains a definition, so
        usage--definition similarity improves only to the extent that the
        usage space improves.
\end{itemize}
